\section{Sistematización del problema y abordaje de la solución}

Para articular el problema investigado con los componentes de la propuesta de ingeniería de software, se presenta la siguiente estructura de correspondencia funcional:

\begin{table}[H]
\centering
\small
\renewcommand{\arraystretch}{1.3}
\begin{tabular}{|p{4cm}|p{4.5cm}|p{4.5cm}|}
\hline
\textbf{Situaciones preliminares bajo investigación (Diagnóstico)} & \textbf{Problema Central y Componente de Solución Propuesta} & \textbf{Consecuencias e Impactos Esperados} \\
\hline
C1: Posible uso predominante de registros manuales o canales de mensajería dispersos para documentar la recepción y los avances técnicos. & A1: Módulo de recepción digitalizada con registro de evidencias físicas (fotografías, observaciones iniciales) y emisión de comprobante único. & F1: Facilitar la verificación transparente del estado de entrega del equipo y brindar respaldo documental claro al cliente y al personal técnico. \\
\hline
C2: Dificultad del cliente para consultar el avance actualizado del mantenimiento sin requerir la interrupción del trabajo del técnico. & A2: Módulo de trazabilidad basado en estados secuenciales del servicio (recepción, diagnóstico, reparación, pruebas, entrega), marcas temporales (timestamps), registro auditable de cambios y restricciones de eliminación en base de datos. & F2: Permitir la consulta organizada de las etapas de atención con registro auditable de cambios, reduciendo la incertidumbre del usuario y las consultas repetitivas. \\
\hline
C3: Demanda de atención informativa repetitiva sobre el funcionamiento del servicio, presupuestos o interpretación de estados de órdenes de trabajo. & A3: Asistente conversacional con Inteligencia Artificial enfocado en responder preguntas frecuentes y consultar estados mediante código de orden verificado. & F3: Agilizar la atención informativa al cliente de manera automatizada, optimizando el tiempo del personal técnico para sus labores operativas. \\
\hline
\end{tabular}
\end{table}